\section{Reading the seed forms on decorated lattices}

% BEGIN B77R1
\paragraph{B77R1. Normalizing the terminal coefficient pairing.}
Fix stages $X$, their original column data $I$, and the resulting
datum $D=X.\mathrm{toD7}$. Fix its family at infinity and the
specified discrete-series model $M_X$ at the terminal space $X'$.
Use the source coefficient construction for $X'$: it gives an
irreducible complex coefficient model $C_X$, its semilinear
coefficient map $c_{0,R}$, the allowed signed Gram matrix and a
nonzero equivariant dual pairing $p$ with the smooth tensor of
$M_X$. At $V_i=V_c(D,i)$ take the positive-definite coefficient
model $C_i$ and its map $c_{0,i}$. These constructions include
the rational action and inner-product identities for their indicated
weights. We can make $p$ isometric in the following precise sense.

Choose an orthonormal basis at each infinite place, writing $b_f$
for the products in the tensor basis. Set
\[
 H(w,w')=\sum_f\overline{p(w)(b_f)}\,p(w')(b_f).
\]
This is invariant under unitary change of each local basis by
orthogonality of the change-of-basis matrices. It is represented
by an endomorphism $T$ of $C_X$, with
$H(w,w')=\langle w,Tw'\rangle$. At each real component the
coefficient action and the dual action on the smooth space
preserve the pairing. The finite factor acts trivially on the
coefficient module. The real and finite factors generate the
adelic action, so $T$ commutes with it. A nonzero eigenspace of
$T$ is invariant and is all of $C_X$ by irreducibility. Thus
$T=r\,\mathrm{id}$. Choose $w$ with $p(w)\ne0$. The functional is nonzero
on some tensor basis vector and
$r=H(w,w)/\|w\|^2>0$. Replace $p$ by $r^{-1/2}p$. The displayed
sum for the new pairing equals $\langle w,w'\rangle$ in every
product orthonormal basis. A nonzero scalar preserves the
dual-pairing covariance. This is exactly $\mathrm{PrIso}$.
In its general source definition the identity is quantified over
finite orthonormal bases. If some local smooth space admits no
such basis, the condition has no instances and the scalar $1$
suffices. No scaling of either $c_{0,R}$ or $c_{0,i}$ is involved.
% END B77R1

% BEGIN B77R2
\paragraph{B77R2. The ordinary row is evaluated on its actual finite matrix.}
Use the original field witness, set of used twists, characters,
group data and trace in $\mathrm{famR}$, and let $L$ be its level.
Assume the column is valid, its character values satisfy
$\mathrm{ValsIn}$, its level satisfies $\mathrm{LevOf}$, and the
specified finite-group identification acts place by place as the
identity. Use the original standard lattices at $X'$ and the $V_i$
and identity level frames $\beta_X,\beta_{V_i}$ as in the seed.
Put $n_R=D.R.N(D.R.m)-1$. For an adelic finite component $h$ of
$X'$, its matrix in the bridge is
\[
 m(h)=\mathrm{bc}(\mathrm{swapAE},h)\,\beta_X.
\]
Here $\mathrm{swapAE}$ puts the coefficient field before the
finite adeles in the tensor product. Two matrices of possibly
different displayed sizes satisfy $\mathrm{MatEq}$ when the sizes
agree and entries at all indices of the same numerical value agree.
For fixed size equality this is precisely equality after reindexing:
extensionality of matrices proves it one entry at a time.

The equality $\operatorname{rank}(X')=n_R$ and the ordinary
weight equality induce the source isomorphism of the rational
coefficient modules. Compose it with $c_{0,R}$ to obtain $c_R$,
before choosing a twist. It takes the standard monomial basis
to the old basis up to its index bijection, so its images are
linearly independent over $\mathbb C$ and span. For a rational
group element $\gamma$, the two reindexed rational matrices
define the same polynomial action, and the rational-action law of
$c_{0,R}$ therefore proves the covariance of $c_R$. The
weight identification carries the two coefficient lattice
operations and the coefficient pairing with it. The source's
signed-Gram comparison on $X'$ gives the pairing attached to
the identity Gram form.

For each twist $t$, use the decorated class presentation on the
marked lattices for $\mathrm{rowSG}(t)$ with empty first prime set.
The standard unitary rational and adelic groups identify with the
groups of $X'$ on these matrices. Decorated points are attained
using the unimodular genus and surjectivity of reduction at the
split marked places. Pull the level character back to the stabilizer.
It is trivial on the principal congruence kernel, and its inverse
is the embedded character of the decorated coefficient action.
Thus a rational move by $\gamma$ and a stabilizer move by $k$
on a decoration act on its coefficient value through
$\chi(k)^{-1}\rho(\gamma)$ after $c_R$. This is the descent law
for the class presentation. Its extension map evaluates the
given decorated form on the decoration of the adelic argument
and applies $c_R$. Restriction followed by its inverse recovers
the original equivariant form on decorations.

To check the universal reading clause, take \emph{any} compatible
CM datum and any unitary $x$ with $\mathrm{MatEq}(m(h),x)$.
Since $\beta_X=1$, its underlying matrix is the reindexed
$\mathrm{bc}(\mathrm{swapAE},h)$, denoted $\mathrm{ptE}(h)$.
Expanding the decoration of the quotient class of $x$ now gives
the very decorated lattice $\mathrm{decoAd}(t,h)$ used by the
presentation. Hence its evaluation is
$c_R(\phi(\mathrm{decoPt}([x])))$ for every source form $\phi$,
where tuple and level on $\mathrm{decoPt}$ are those just fixed.
For existence of such $x$ take the stages' own CM datum once:
$h$ is the finite component of an adelic unitary element and the
terminal Gram matrix is one, so $\mathrm{ptE}(h)$ is unitary for
the standard form. Taking this element witnesses $\mathrm{MatEq}$
for every $h$. This gives all four clauses of $\mathrm{BrgStdR}$.
% END B77R2

% BEGIN B77R3
\paragraph{B77R3. The lift uses one coefficient map for every twist of a row.}
Keep B77R2's data and fix a lift row $i$. Here the decorated rank
is $n_V$, and the weight is the specified $\lambda_V(i+1)$.
Transport along $n_V=\operatorname{rank}(V_i)$ and the equality
of this weight with the model's forced lift weight. The resulting
rational isomorphism $e_i$ intertwines the $\rho_{V,r}$ actions
and transports the basis, pairing and both distinguished lattice
operations. Put $c_i=c_{0,i}e_i$ before choosing $t$.
For a rational matrix $\gamma_E$ representing $\gamma$, the
intertwining identity applied to its reindexing, then that of
$c_{0,i}$, gives $c_i\rho_{V,r}(\gamma_E)=\rho(\gamma)c_i$ in
these coordinates. A witness of $\mathrm{MatEq}$ for any alleged
representative identifies it with this reindexing.

The base sign is positive at $i+1$. Thus the signed Gram form
in the lift coefficient pairing is the original Gram form, which
is one on this row. In particular the identity rational frame
and identity finite frame meet the unitary frame requirements.
For each $t$ use the level and character in these frames. As on
the ordinary side, the decorated-point map is onto and the
level character is trivial on the principal kernel. For the
rational and stabilizer pair $(\gamma,k)$ the action on a
decorated coefficient is the rational matrix action times the
embedded inverse character at $k$. The two successive rank
reindexings of the matrix of $\gamma$ cancel. Applying $c_i$
gives $\chi(k)^{-1}\rho(\gamma)c_i$, the required equivariance.
Together with the transported inner product and basis this
constructs the class presentation $P_{i,t}$. Its coefficient
map is $c_i$ for all $t$. Denote its extension by $j_{i,t}$.

Now take $h$ in the finite adelic group of $V_i$ and any admissible
$x$ in the standard unitary model whose matrix is
$\mathrm{MatEq}$ to $\mathrm{bc}(\mathrm{swapAE},h)\beta_{V_i}$.
The identity frame gives $x$ exactly the reindexed finite matrix.
Expansion of the prescribed decorated point at $(1,1)$ gives the
decoration of $[x]$. The presentation therefore gives, on a
restricted equivariant decorated form,
\[
 j_{i,t}(\mathrm{res}\,\phi)(h)
           = c_i\bigl(\phi(\mathrm{decoPt}([x]))\bigr).
\]
The basis images in the base tuple presentation are linearly
independent. These are the images under the same $c_i$. The
original CM datum works for every $h$. Its finite component preserves the true Gram form, hence the
identity form, and transporting both indices carries that equation
to rank $n_V$. This is a unitary $x$ of exactly the required matrix.
We obtain $\mathrm{BrgStdV}$ at row $i$. In both readings one
may now fix the CM datum and choose a matrix representative for
$h$ first. For each $t$ its decorated class works for all forms
$\phi$ simultaneously.
% END B77R3

% BEGIN B77R4
\paragraph{B77R4. What a double-coset comparison gives on these readings.}
Fix an allowed split prime $P$ and degree $k$. The source operator
on decorated forms sums over the finite set of its $(P,k)$
neighbors. In degree zero the neighbor operation is the identity.
In degree $k>n$ there are no such neighbors for a space of decorated
rank $n$: their count is the Gaussian binomial for $k$-dimensional
subspaces of a residue space of dimension $n$, hence zero. The sum
is then the zero operator. Restriction transports both identities
to the class presentation. Use $n=n_R$ in B77R2 and $n=n_V$ in
B77R3. These are simultaneous statements in all twists and rows.

For $1\le k\le n$, let $A_{P,k}$ be the finite set of representatives
in the original split local comparison. Its substantive local inputs
are that it represents the specified double coset, consists of
elements supported at the permitted place, and, for every $g$, that
$a\mapsto d(ga)$ bijects $A_{P,k}$ with the decorated neighbors
of $d(g)$. Here $d(g)$ is the decoration in the very presentation
being read, with the same level. On the ordinary row let
$\phi$ denote the equivariant form underlying a class form $f$.
B77R2 and additivity of $c_R$ give
\[
  j_R(T_{P,k} f)(g)
   =\sum_{d'\text{ neighbor of }d(g)} c_R(\phi(d'))
   =\sum_{a\in A_{P,k}} j_R(f)(ga).
\]
The last equality reindexes by that bijection. In particular the
comparison uses the whole finite sum, with no division by its
cardinality. On the lift side the group isomorphism given by the
identity frames sends the representatives to the lift group and
carries the standard local compact subgroup to the intersection
of the level group with the local support subgroup. The source
local decomposition extends this double coset to the level
group. Transport its neighbor bijection and apply B77R3 term by
term, now with $c_i$. This gives the same operator identity on
the actual lift presentation.

For the full seed bridge the action inputs also specify the
reduction of each level element by its flag tuple. The inert
operator inputs select the original list of elementary divisors,
its allowed label, its support and the corresponding finite
double-coset representatives on these same presentations.
The split neighbor bijections, these action and inert comparisons,
and the original place, root and character value conditions are
retained here as source inputs. Insert the ordinary presentations
and the lift presentations into one $B$ with ranks $r+mt$ and
$r+t$, respectively, where $t$ in these two ranks is the rank
increment. The split-place map and the inert-place map are the
respective original inclusions. The split bad set is $L.P_g$,
and the finite set on the real field consists of its contractions.
The covariance, coefficient ties and two independent degree
bounds above are on the maps $B.j_R,B.j_L$. This identifies the
adelic readings used in the seed construction with the
decorated-lattice evaluations of B77R2 and B77R3.
% END B77R4

% BEGIN B77R5
\paragraph{B77R5. A single set of representatives for the split neighbors.}
We supply the decorated-neighbor bijection in B77R4. Work first in
the standard unitary model of rank $n$. Let $P$ be an allowed split
prime, $Q=\bar P$ its distinct conjugate, and $v$ the real prime
below them. Neither prime belongs to the marked decoration sets.
Let $K_v$ be the group supported at $v$ stabilizing the standard
local lattice. Choose $\pi$ with valuation one at $Q$ and zero
at $P$. Let $\eta$ be supported at $v$, with diagonal
$\pi/\bar\pi$ on the first $k$ coordinates and one on the others,
for $1\le k\le n$. Its conjugate times itself is one. In local
lattice coordinates the first $k$ directions therefore acquire
exponent $+1$ at $Q$ and exponent $-1$ at $P$. The other directions
and all other places have exponent zero.

Let $g$ be any standard finite unitary element, with lattice $N$.
In its frame, reducing all $n$ coordinates at $Q$ gives a
surjection $r:N\to k(Q)^n$: localization gives the full local
lattice, and denominators prime to $Q$ can be inverted modulo $Q$.
For $g\eta$ the intersection of its lattice with $N$ is the kernel
of the first $k$ coordinates of $r$. The contraction at $Q$ gives
these $k$ conditions, and the enlargement at $P$ adds none. The
neighbor is self-dual and has the same decorations away from $v$.
Left multiplication of $\eta$ by an element of $K_v$ changes the
residue frame, and right multiplication by $K_v$ changes neither
the resulting lattice nor its decorations.

Conversely take a decorated $P$-neighbor $N'$ of $N$. Put
$M=N\cap N'$. The neighbor conditions give $Q N\subseteq M$
and $N/M\simeq k(Q)^k$. Thus the quotient map factors through
$r$ as a surjection $k(Q)^n\to k(Q)^k$. Its $k$ coordinate
linear forms are independent. Complete them to the rows of an
invertible $n$ by $n$ residue matrix. Lift this matrix to the
integral local stabilizer: an integral entrywise lift has unit
determinant, and the inverse conjugate transpose prescribes
the other factor of the split unitary group. Extend it by the
identity away from $v$. Use its inverse as the frame change
$\kappa\in K_v$, since pullback of coordinates uses the inverse
matrix. The lattice of $g\kappa\eta$ now has intersection $M$
with $N$. The self-dual gluing formula for a $P$-neighbor is
\[
                    N'=M+Q M^\vee .
\]
At $Q$ one has $M_Q=N'_Q$ and
$(QM^\vee)_Q=Q N_Q\subseteq M_Q$. At $P$, the ideal $Q$ is a
unit and the dual of $M_Q=N'_Q$ in the paired completion is $N'_P$.
At every other prime all three lattices $N,N',M$ agree and are
locally self-dual across the conjugate pair. The formula follows
by equality at every localization. Applying this to both neighbors
gives equality of their full lattices. The flags and residue bases
at marked primes also agree, because the move is supported away
from those primes. This proves surjectivity on decorated neighbors.

If two elements of $K_v\eta K_v$ give the same decoration of $g$,
their global lattices coincide. Localizing, canceling $g$, and
using that the two elements are the identity off $v$, shows that
their quotient stabilizes the standard local lattice and lies in
$K_v$. Thus the two cosets $aK_v$ coincide. At $g=1$ the resulting
bijection identifies the cosets $aK_v$ with a finite neighbor set.
Choose one representative of each once. For \emph{every} $g$ the
same finite set enumerates the neighbors without repetition.
Transport it by the unitary group identification of B77R2 or
B77R3. That identification respects the actual finite matrices,
so the image of $g a$ gives the same decoration on the reading
side, as required for B77R4. This supplies the split neighbor
bijection. The additional determinant, support and full-level
double-coset identities in the operator tie and the inert
comparisons have their separate inputs in B77R4.
% END B77R5
